\section{Introduction}
\label{sec:introduction}

Scientific research is an iterative and knowledge-intensive process in
which researchers continuously generate ideas, formulate hypotheses,
conduct experiments, analyze results, and produce scientific artifacts.
During this process, valuable knowledge is distributed across multiple
sources including papers, laboratory notes, software repositories,
discussion records, and project documents. The increasing complexity of
modern research projects has made maintaining continuity, traceability,
and reproducibility a significant challenge.

Traditional research management systems mainly focus on document
organization, workflow execution, or project scheduling. Although these
approaches provide valuable support for storing and accessing
information, they generally do not preserve the evolution of research
reasoning itself. Important elements such as why a decision was made,
which hypothesis led to an experiment, how results influenced future
directions, and how artifacts were derived are often lost or remain
implicitly stored in researchers' personal memory.

Recent advances in knowledge graphs, semantic retrieval, provenance
models, and generative artificial intelligence have created new
opportunities for intelligent research assistance. However, existing
approaches typically address isolated aspects of knowledge management.
Workflow systems focus on computational execution, semantic systems
emphasize information representation, and retrieval-based approaches
mainly improve access to existing documents. A unified framework for
maintaining persistent research memory across the complete research
lifecycle remains insufficiently explored.

This work introduces the Cognitive Research Memory Engine (CRME), a
structured framework designed to address the problem of research
knowledge continuity. Instead of treating research information as
independent documents, CRME represents research activities as connected
memory objects containing ideas, hypotheses, decisions, experiments,
results, artifacts, and publication-related information. These objects
form a persistent research memory that enables historical reconstruction,
semantic organization, and provenance-aware retrieval.

The main objective of CRME is not to replace researchers or automate
scientific discovery completely, but to provide a cognitive memory layer
that supports researchers during long-term projects. By maintaining
structured records of research evolution, CRME aims to reduce knowledge
loss, improve reproducibility, and support more explainable research
workflows.


The main contributions of this paper are summarized as follows:

\begin{itemize}

\item A structured research memory model is introduced, representing
research activities as interconnected knowledge objects across the
research lifecycle.

\item A modular architecture for persistent research memory management is
presented, integrating memory storage, session continuity, semantic
organization, provenance tracking, and project evolution.

\item A provenance-aware approach is proposed to preserve relationships
between research decisions, experiments, results, and generated
artifacts.

\item An implementation-oriented evaluation methodology is defined,
including capability comparison, component validation, and reproducibility
considerations.

\end{itemize}


The remainder of this paper is organized as follows.
Section~\ref{sec:related_work} discusses related research management,
knowledge representation, provenance, and retrieval approaches.
Section~\ref{sec:methodology} presents the CRME architecture and internal
research memory model.
Section~\ref{sec:evaluation} describes the evaluation methodology and
validation strategy.
Finally, conclusions and future research directions are discussed.



